% GENERATED FILE. Edit canonical lessons, then run npm run generate:books.
% canonical-source-hash: fnv1a64:f81747c57527d10e
% canonical-lessons: SA-C360-yasah, SA-C360-pratistha, SA-C360-ahankarah, SA-C360-apamanah, SA-C360-pratarayati

\chapter{Honour and Dishonour}
\label{ch:sa-honour-and-dishonour}
\hlchaptermodality{360}

\begin{chapteropening}
\textbf{By the end of this chapter:} \emph{I can talk about fame, standing and insults.}

Complete the last lesson of chapter 360: I can talk about fame, standing and insults.
\end{chapteropening}

\section[\texorpdfstring{\sk{यशः} (yaśaḥ)}{यशः (yaśaḥ)}]{\sk{यशः} (yaśaḥ) — fame; glory}

\label{lesson:SA-C360-yasah}



\begin{quote}
Before the new one: say the Sanskrit for craving, then the Sanskrit for dedicates.
\end{quote}

\subsection*{You'll want to know: \sk{यशः}}
\textbf{\sk{यशः}} — \emph{yaśaḥ} — \textquotedblleft{}fame; glory\textquotedblright{}.

\begin{quote}
\sk{लेखकस्य} \sk{यशः} \sk{सर्वत्र} \sk{अस्ति।}
\end{quote}

\begin{quote}
\emph{lekhakasya yaśaḥ sarvatra asti.}
\end{quote}

\begin{quote}
The writer's fame is everywhere.
\end{quote}

A neuter noun.

\begin{cousinweb}
\textbf{\sk{यशः}} (\emph{yaśaḥ}) MW gives \emph{yaśas} \textquotedblleft{}honour, glory, fame, renown\textquotedblright{} and gives no derivation, so its origin is uncertain (MW).
\end{cousinweb}

\begin{grammarlens}[title={What you've built}]
One more word for giving reasons and opinions.
\end{grammarlens}

\subsection*{Your turn}
\begin{itemize}
\raggedright
  \item \emph{Say it:} \emph{yaśaḥ}
  \item \emph{Say it:} \emph{yaśaḥ}, once more
  \item \emph{Say it:} \emph{lekhakasya yaśaḥ sarvatra asti.} — The writer's fame is everywhere.
  \item \emph{Recall:} say the Sanskrit for craving, then the Sanskrit for dedicates, then say \emph{yaśaḥ} again
\end{itemize}

\begin{tcolorbox}[breakable,colback=teal!4,colframe=teal!35!black,title={Before you move on}]
What does \sk{यशः} mean? (\textquotedblleft{}Fame; glory\textquotedblright{}.) Say it once more.
\end{tcolorbox}
\section[\texorpdfstring{\sk{प्रतिष्ठा} (pratiṣṭhā)}{प्रतिष्ठा (pratiṣṭhā)}]{\sk{प्रतिष्ठा} (pratiṣṭhā) — standing; reputation}

\label{lesson:SA-C360-pratistha}



\begin{quote}
Before the new one: say the Sanskrit for dedicates, then the Sanskrit for glory.
\end{quote}

\subsection*{You'll want to know: \sk{प्रतिष्ठा}}
\textbf{\sk{प्रतिष्ठा}} — \emph{pratiṣṭhā} — \textquotedblleft{}standing; reputation\textquotedblright{}.

\begin{quote}
\sk{वैद्यस्य} \sk{ग्रामे} \sk{प्रतिष्ठा} \sk{अस्ति।}
\end{quote}

\begin{quote}
\emph{vaidyasya grāme pratiṣṭhā asti.}
\end{quote}

\begin{quote}
The doctor has standing in the village.
\end{quote}

A feminine noun.

\begin{cousinweb}
\textbf{\sk{प्रतिष्ठा}} (\emph{pratiṣṭhā}) is built on \emph{prati} and the root √sthā, \textquotedblleft{}to stand\textquotedblright{}; MW gives \emph{prati-ṣṭhā} \textquotedblleft{}steadfastness\textquotedblright{} and \textquotedblleft{}high rank or position, fame\textquotedblright{} (MW).
\end{cousinweb}

\begin{grammarlens}[title={What you've built}]
One more word for giving reasons and opinions.
\end{grammarlens}

\subsection*{Your turn}
\begin{itemize}
\raggedright
  \item \emph{Say it:} \emph{pratiṣṭhā}
  \item \emph{Say it:} \emph{pratiṣṭhā}, once more
  \item \emph{Say it:} \emph{vaidyasya grāme pratiṣṭhā asti.} — The doctor has standing in the village.
  \item \emph{Recall:} say the Sanskrit for dedicates, then the Sanskrit for glory, then say \emph{pratiṣṭhā} again
\end{itemize}

\begin{tcolorbox}[breakable,colback=teal!4,colframe=teal!35!black,title={Before you move on}]
What does \sk{प्रतिष्ठा} mean? (\textquotedblleft{}Standing; reputation\textquotedblright{}.) Say it once more.
\end{tcolorbox}
\section[\texorpdfstring{\sk{अहंकारः} (ahaṅkāraḥ)}{अहंकारः (ahaṅkāraḥ)}]{\sk{अहंकारः} (ahaṅkāraḥ) — arrogance; ego}

\label{lesson:SA-C360-ahankarah}



\begin{quote}
Before the new one: say the Sanskrit for glory, then the Sanskrit for standing.
\end{quote}

\subsection*{You'll want to know: \sk{अहंकारः}}
\textbf{\sk{अहंकारः}} — \emph{ahaṅkāraḥ} — \textquotedblleft{}arrogance; ego\textquotedblright{}.

\begin{quote}
\sk{अहंकारः} \sk{मित्राणि} \sk{दूरं} \sk{नयति।}
\end{quote}

\begin{quote}
\emph{ahaṅkāraḥ mitrāṇi dūraṃ nayati.}
\end{quote}

\begin{quote}
Arrogance drives friends away.
\end{quote}

A masculine noun.

\begin{cousinweb}
\textbf{\sk{अहंकारः}} (\emph{ahaṅkāraḥ}) is built from \emph{aham}, \textquotedblleft{}I\textquotedblright{}, and \emph{kāra}, \textquotedblleft{}a making\textquotedblright{}; MW gives \emph{ahaṃ-kāra} \textquotedblleft{}the making of self, egotism\textquotedblright{} and \textquotedblleft{}pride, haughtiness\textquotedblright{} (MW).
\end{cousinweb}

\begin{grammarlens}[title={What you've built}]
One more word for giving reasons and opinions.
\end{grammarlens}

\subsection*{Your turn}
\begin{itemize}
\raggedright
  \item \emph{Say it:} \emph{ahaṅkāraḥ}
  \item \emph{Say it:} \emph{ahaṅkāraḥ}, once more
  \item \emph{Say it:} \emph{ahaṅkāraḥ mitrāṇi dūraṃ nayati.} — Arrogance drives friends away.
  \item \emph{Recall:} say the Sanskrit for glory, then the Sanskrit for standing, then say \emph{ahaṅkāraḥ} again
\end{itemize}

\begin{tcolorbox}[breakable,colback=teal!4,colframe=teal!35!black,title={Before you move on}]
What does \sk{अहंकारः} mean? (\textquotedblleft{}Arrogance; ego\textquotedblright{}.) Say it once more.
\end{tcolorbox}
\section[\texorpdfstring{\sk{अपमानः} (apamānaḥ)}{अपमानः (apamānaḥ)}]{\sk{अपमानः} (apamānaḥ) — an insult; dishonour}

\label{lesson:SA-C360-apamanah}



\begin{quote}
Before the new one: say the Sanskrit for standing, then the Sanskrit for arrogance.
\end{quote}

\subsection*{You'll want to know: \sk{अपमानः}}
\textbf{\sk{अपमानः}} — \emph{apamānaḥ} — \textquotedblleft{}an insult; dishonour\textquotedblright{}.

\begin{quote}
\sk{सः} \sk{अपमानं} \sk{सहते}, \sk{किन्तु} \sk{न} \sk{विस्मरति।}
\end{quote}

\begin{quote}
\emph{saḥ apamānaṃ sahate, kintu na vismarati.}
\end{quote}

\begin{quote}
He bears an insult but does not forget it.
\end{quote}

A masculine noun.

\begin{cousinweb}
\textbf{\sk{अपमानः}} (\emph{apamānaḥ}) is built on \emph{apa}, \textquotedblleft{}away\textquotedblright{}, and \emph{māna}, \textquotedblleft{}honour\textquotedblright{}, from the root √man; MW gives \emph{apa-māna} \textquotedblleft{}disrespect, contempt, disgrace\textquotedblright{} (MW).
\end{cousinweb}

\begin{grammarlens}[title={What you've built}]
One more word for giving reasons and opinions.
\end{grammarlens}

\subsection*{Your turn}
\begin{itemize}
\raggedright
  \item \emph{Say it:} \emph{apamānaḥ}
  \item \emph{Say it:} \emph{apamānaḥ}, once more
  \item \emph{Say it:} \emph{saḥ apamānaṃ sahate, kintu na vismarati.} — He bears an insult but does not forget it.
  \item \emph{Recall:} say the Sanskrit for standing, then the Sanskrit for arrogance, then say \emph{apamānaḥ} again
\end{itemize}

\begin{tcolorbox}[breakable,colback=teal!4,colframe=teal!35!black,title={Before you move on}]
What does \sk{अपमानः} mean? (\textquotedblleft{}An insult; dishonour\textquotedblright{}.) Say it once more.
\end{tcolorbox}
\section[\texorpdfstring{\sk{प्रतारयति} (pratārayati)}{प्रतारयति (pratārayati)}]{\sk{प्रतारयति} (pratārayati) — misleads; deceives}

\label{lesson:SA-C360-pratarayati}



\begin{quote}
Before the new one: say the Sanskrit for arrogance, then the Sanskrit for an insult.
\end{quote}

\subsection*{You'll want to know: \sk{प्रतारयति}}
\textbf{\sk{प्रतारयति}} — \emph{pratārayati} — \textquotedblleft{}misleads; deceives\textquotedblright{}.

\begin{quote}
\sk{धूर्तः} \sk{जनान्} \sk{प्रतारयति।}
\end{quote}

\begin{quote}
\emph{dhūrtaḥ janān pratārayati.}
\end{quote}

\begin{quote}
The rogue misleads people.
\end{quote}

Past (he …): \sk{सः} \sk{मित्रं} \sk{प्रतारितवान्}, \emph{saḥ mitraṃ pratāritavān}.

\begin{cousinweb}
\textbf{\sk{प्रतारयति}} (\emph{pratārayati}) is built on \emph{pra} and the root √tṝ, \textquotedblleft{}to cross\textquotedblright{}, in its causative form; MW gives \emph{pra-√tṝ} \textquotedblleft{}to mislead, to take in, to deceive\textquotedblright{} and \emph{pra-tārita} \textquotedblleft{}misled, deceived\textquotedblright{} (MW).
\end{cousinweb}

\begin{grammarlens}[title={What you've built}]
One more word for giving reasons and opinions.
\end{grammarlens}

\subsection*{Your turn}
\begin{itemize}
\raggedright
  \item \emph{Say it:} \emph{pratārayati}
  \item \emph{Say it:} \emph{pratārayati}, once more
  \item \emph{Say it:} \emph{dhūrtaḥ janān pratārayati.} — The rogue misleads people.
  \item \emph{Recall:} say the Sanskrit for arrogance, then the Sanskrit for an insult, then say \emph{pratārayati} again
\end{itemize}

\begin{tcolorbox}[breakable,colback=teal!4,colframe=teal!35!black,title={Before you move on}]
What does \sk{प्रतारयति} mean? (\textquotedblleft{}Misleads; deceives\textquotedblright{}.) Say it once more.
\end{tcolorbox}
